\begin{remark}[E12: what a tag-value filter accepts]\label{rem:glob-constructs}
Definition~\ref{def:glob} is the premise of every negative result in
\S\ref{sec:regular}, and until the eleventh pass it was a sentence from Datadog's
documentation. That is the shape of error this paper's method note is about: a rule read
in a manual and never put to the system. It is also load-bearing in the reverse
direction --- Table~\ref{tab:regex-glob-reasons} refuses $140$ matcher occurrences on
the ground that no glob denotes a single arbitrary character, which is $13$ points of
the $37.3\%$ the abstract quotes --- so if the documentation were incomplete, that row
would be the paper's largest single error.

Three series were submitted to the trial account, tagged \texttt{probe:web-1},
\texttt{probe:web-12} and \texttt{probe:web-a}, with a control query carrying no glob to
establish that all three were stored, which is the discipline E5 added. Then three
questions, each designed so that the construct's presence and absence give different
answers:

\begin{center}\footnotesize
\begin{tabular}{@{}lll@{}}
\toprule
Construct & Scope asked & If live \\
\midrule
single-character wildcard & \texttt{probe:web-?} & selects \texttt{web-1}, \texttt{web-a} \\
character class & \texttt{probe:web-[0-9]} & selects \texttt{web-1} \\
alternation in a value & \texttt{probe:web-1|web-a} & selects two \\
\bottomrule
\end{tabular}
\end{center}

\noindent
The control stored all three series. All three questions came back as \emph{parse
errors} --- \texttt{Rule 'or\_op' didn't match at '?'}, and the same at \texttt{[} and
at \texttt{|}. That is a better answer than the one the experiment was designed to get.
A construct that is accepted and selects nothing would be a literal, and a literal is
still a fact about matching that a future Datadog release could change quietly; a
construct the query grammar refuses to parse is absent from the language. Definition~14
is now a measurement, and the $140$ occurrences rest on it rather than on
\citet{ddfilter}.

Two limits worth stating. This is one account on one site on one date, and Datadog
versions its query grammar silently, so the reading is of the service as it behaved
rather than of a specification. And a parse error is evidence about the \emph{filter}
grammar only: it says nothing about what the editor's autocomplete offers or what a
different API surface accepts, both of which are outside what any query in this corpus
exercises.

This is the ninth vendor measurement and the sixth to confirm rather than correct. The
count is worth keeping visible in both directions: the three that corrected something
are the reason the other six were worth running, and six confirmations in a row would
be a reason to wonder whether the questions had got too easy.
\end{remark}
